\documentclass[10pt,a4paper]{amsart}
\usepackage[margin=19mm]{geometry}
\usepackage{amsmath,amssymb,mathtools}
\usepackage[scr=rsfs,cal=euler]{mathalfa}
\usepackage{xfrac}
\usepackage[unicode,hidelinks,bookmarks=false]{hyperref}
\newcommand{\ie}{i.e.\ }
\newcommand{\cf}{cf.\ }
\newcommand{\resp}{respectively\ }
\newcommand{\Subsection}{\subsection}
\providecommand{\nobreakdash}{\nobreak\hbox{-}}
\setlength{\emergencystretch}{2em}
\multlinegap0pt
\usepackage{amssymb}
\usepackage{mathtools}
\usepackage{csquotes}
\usepackage{float}
\newtheorem{proposition}{Proposition}
\title{Signature Kernel and Schwinger-Dyson Kernel Equations as Two-Parameter Rough Differential Equations}
\author{Thomas Cass, Dan Crisan, Andrea Iannucci, William F. Turner}
\date{Source: arXiv:2605.08844v1 (CC BY 4.0)}
\begin{document}
\maketitle

\noindent\emph{Source and licence.} Signature Kernel and Schwinger-Dyson Kernel Equations as Two-Parameter Rough Differential Equations. Version arXiv:2605.08844v1 (CC BY 4.0), distributed by arXiv under the Creative Commons Attribution 4.0 International licence, http://creativecommons.org/licenses/by/4.0/, as recorded in the arXiv licence metadata for this exact version and shown on the abstract page https://arxiv.org/abs/2605.08844v1. Full licence notice: \texttt{licenses/arxiv-2605.08844-CC-BY-4.0.txt}.\par\smallskip

\noindent\textit{Editorial scope.} Counted unit: the article's proposition stability\_proposition (source0.tex lines 1092--1103). The article's complete original proof of it is delivered with it (source0.tex lines 2477--2683, one \texttt{proof} environment of 207 lines, which the article places away from the statement). No mathematics is written, repaired or re-derived: the statement and the proof are the article's own lines, in the article's own notation, and no retained line is substituted or re-flowed.  Retained uncounted as the source-local context the proof needs: lines 535--540.  Nothing was omitted from any retained span, so the package carries no omission record.  No retained line contains a citation, so the wrapper carries no bibliography and no bibliography entry.  No retained line contains a figure environment, an image-inclusion command or drawing code, so no image is delivered and the package declares no asset dependency.  Editorial formatting only, in addition: an AMS \texttt{amsart} preamble that restates the article's own declarations and macros used by the retained lines, namely the environments \texttt{proposition} on separate counters, together with the standard packages those lines need.  The retained environments print in this document as Proposition 1 in order of appearance, whereas the article prints the same items with the numbering of its own full document.  The complete native source of the article as distributed in the arXiv e-print for this exact version is bundled unchanged as \texttt{sources/200178/source0.tex}.
\begin{equation}\label{remainder_remainder_norm}
\left\| \mathbf{R} \begin{pmatrix} s,t \\ u,v \end{pmatrix} [\eta_j \boxtimes \xi_k] \right\| \leq C \omega_{\mathbf{X}}\left(
s,t\right)  ^{\left(  \kappa-j\right)  /p} \omega_{\mathbf{Y}}\left(
u,v\right)  ^{\left(  \lambda-k\right)  /q}\left \| \eta_{j}%
\right \| \| \xi_k \|,
\end{equation}

\begin{proposition}\label{stability_proposition}
    Let $\mathbf{X}, \tilde{\mathbf{X}} \in \Omega G_{p, \omega_\mathbf{X}}([a,b], V)$, $\mathbf{Y}, \tilde{\mathbf{Y}} \in \Omega G_{q, \omega_\mathbf{Y}}([c,d], W)$ and  $Z \in \mathcal{D}^{p,q}_{\mathbf{X},\mathbf{Y}}$, $\tilde{Z} \in \mathcal{D}^{p,q}_{\tilde{\mathbf{X}},\tilde{\mathbf{Y}}}$. The composition of $Z$ with $g$ satisfies 
    \[
    \|g(Z)\|_{\mathcal{D}^{p,q}_{\mathbf{X}, \mathbf{Y}}} \leq C_{p,q,g,\omega_\mathbf{X}, \omega_\mathbf{Y}, Z} \|Z\|_{\mathcal{D}^{p,q}_{\mathbf{X}, \mathbf{Y}}},
    \]
    moreover the following stability estimate holds
    \[
    d(g(Z), g(\tilde{Z})) \leq C_{p,q,g,\omega_\mathbf{X}, \omega_\mathbf{Y}, Z, \tilde{Z}} d(Z, \tilde{Z}).
    \]


\end{proposition}

\begin{proof}[Proof of Proposition \ref{stability_proposition}]
To prove the  continuity property we only need to check that $R^{g,1}$ and $R^{g,2}$ are respectively  $\mathbf{Y}$ and $\mathbf{X}$-controlled paths. Since the proof is the analogous for both processes we will only consider the path $R^{g,1}$.\newline
We begin by recalling the decomposition of $R^{g,1}$ derived in the previous proof, in which we are writing
\[
R^{g, 1} = A + B + C + E,
\]
 in order to show that $R^{g,1}$  is controlled by $\mathbf{Y}$ it is sufficient to show that each term in the above sum is controlled by $\mathbf{Y}$.\newline 
Starting with $A$, fix $u$ and expand the factors in the $v$--variable around $u$. This yields
\begin{align*}
A_{s,t}(v)
&=
\sum_{\ell=1}^{\lambda+\kappa-2}\sum_{i=0}^{\ell}\frac{1}{\ell!}\,
D^\ell g\left(Z^{(0,0)}_{t,v}\right)
\circ
\left(
Z^{\otimes(i-1)}_{s,u} \circ [L_{\mathbf{X}^{<\kappa}_{s,t}} \boxtimes id] \otimes R^1_{s,t}(u)\otimes Z^{\otimes(\ell-i)}_{t,u}
\right)
\circ
\left[id \boxtimes L_{\mathbf{Y}^{<\lambda}_{u,v}}\right]^{\otimes\ell}
\circ (\Delta^{\boxtimes}_{>0})^\ell
\\
&\quad + \hat{\mathbf{R}}^{A}\begin{pmatrix} s,t \\ u, v \end{pmatrix},
\end{align*}
where $\hat{\mathbf{R}}^{A}\begin{pmatrix} s,t \\ u, v \end{pmatrix}$ collects the remaining contributions
coming from terms containing one (or more) remainder terms in the expansions of
$Z_{s,v}$, $Z_{t,v}$, and $R^1_{s,t}(v)$.

We now focus on the leading term and reindex the sums as
$i=0,\dots,\lambda+\kappa-2$ and $\ell=i,\dots,\lambda+\kappa-2$:
\begin{equation}\label{pre_Taylor_g_Z_second_rem}
\begin{aligned}
&\sum_{i=0}^{\lambda+\kappa-2}\sum_{\ell=i}^{\lambda+\kappa-2}\frac{1}{\ell!}\,
D^\ell g\left(Z^{(0,0)}_{t,v}\right)
\circ
\left(
 Z^{\otimes(i-1)}_{s,u} \circ [L_{\mathbf{X}^{<\kappa}_{s,t}} \boxtimes id]\otimes R^1_{s,t}(u)\otimes Z^{\otimes(\ell-i)}_{t,u}
\right)
\circ
\left[id \boxtimes L_{\mathbf{Y}^{<\lambda}_{u,v}}\right]^{\otimes\ell}
\circ (\Delta^{\boxtimes}_{>0})^\ell \\
& \qquad +\hat{\mathbf{R}}^{A}\begin{pmatrix} s,t \\ u, v \end{pmatrix}.
\end{aligned}
\end{equation}

Next, we expand $D^\ell g$ around $Z^{(0,0)}_{t,u}$. For each $\ell$ we obtain
\begin{align*}
&D^\ell g\left(Z^{(0,0)}_{t,v}\right)
=
\sum_{m=\ell}^{\lambda+\kappa-2}\frac{1}{(m-\ell)!}\,
D^m g\left(Z^{(0,0)}_{t,u}\right)
\circ
\left(Z^{(0,0)}_{t,v}-Z^{(0,0)}_{t,u}\right)^{\otimes(m-\ell)}
+ \mathrm{Rem}_{\ell}(u,v),
\end{align*}
and we absorb the Taylor remainder $Rem_\ell(u,v)$ into
$\hat{\mathbf{R}}^{A}\begin{pmatrix} s,t \\ u, v \end{pmatrix}$. It immediate to see that this term satisfies the constraint \eqref{remainder_remainder_norm}. Substituting this Taylor in \eqref{pre_Taylor_g_Z_second_rem} yields
\begin{align*}
&=
\sum_{i=0}^{\lambda+\kappa-2}
\sum_{\ell=i}^{\lambda+\kappa-2}
\sum_{m=\ell}^{\lambda+\kappa-2}
\frac{1}{\ell!(m-\ell)!}\,
D^m g\left(Z^{(0,0)}_{t,u}\right)
\\
&\qquad\circ
\left(
Z^{\otimes(i-1)}_{s,u}\otimes R^{1}_{s,t}(u)\otimes Z^{\otimes(\ell-i)}_{t,u}
\otimes
\left(Z^{(0,0)}_{t,v}-Z^{(0,0)}_{t,u}\right)^{\otimes(m-\ell)}
\right)
\circ
\left[id \boxtimes L_{\mathbf{Y}^{<\lambda}_{u,v}}\right]^{\otimes\ell}
\circ (\Delta^{\boxtimes}_{>0})^\ell
\\
&\qquad + \hat{\mathbf{R}}^{A}\begin{pmatrix} s,t \\ u, v \end{pmatrix}.
\end{align*}

Finally, reindex the triple sum by fixing $m$ and summing over $\ell\le m$.
Using $\binom{m}{m-\ell}=\frac{m!}{\ell!(m-\ell)!}$, we obtain
\begin{align*}
&=
\sum_{i=0}^{\lambda+\kappa-2}
\sum_{m=i}^{\lambda+\kappa-2}
\sum_{\ell=i}^{m}
\frac{1}{m!}\binom{m}{m-\ell}\,
D^m g\left(Z^{(0,0)}_{t,u}\right)
\\
&\qquad\circ
\left(
Z^{\otimes(i-1)}_{s,u}\otimes R^{1}_{s,t}(u)\otimes Z^{\otimes(\ell-i)}_{t,u}
\otimes
\left(Z^{(0,0)}_{t,v}-Z^{(0,0)}_{t,u}\right)^{\otimes(m-\ell)}
\right)
\circ
\left[id \boxtimes L_{\mathbf{Y}^{<\lambda}_{u,v}}\right]^{\otimes\ell}
\circ (\Delta^{\boxtimes}_{>0})^\ell
\\
&\qquad + \hat{\mathbf{R}}^{A}\begin{pmatrix} s,t \\ u, v \end{pmatrix}.
\end{align*}
We can now expand the term $Z^{(0,0)}_{t,v} - Z^{(0,0)}_{t,u}$ and include the terms that present at least one remainder in $\hat{\mathbf{R}}^{A}\begin{pmatrix} s,t \\ u, v \end{pmatrix}$.\newline
\begin{align*}
&=
\sum_{i=0}^{\lambda+\kappa-2}
\sum_{m=i}^{\lambda+\kappa-2}
\sum_{\ell=i}^{m}
\frac{1}{m!}\binom{m}{m-\ell}\,
D^m g\left(Z^{(0,0)}_{t,u}\right)
\\
&\qquad \circ \left(
\left(
Z^{\otimes(i-1)}_{s,u}\otimes R^{1}_{s,t}(u)\otimes Z^{\otimes(\ell-i)}_{t,u}
\right)
\circ 
\left[id \boxtimes L_{\mathbf{Y}^{<\lambda}_{u,v}}\right]^{\otimes\ell}
\circ (\Delta^{\boxtimes}_{>0})^\ell \otimes
\left(Z_{t,u} \circ [ \mathbf{1} \boxtimes \mathbf{Y}^{\geq 1}_{s,t}  ]\right)^{\otimes(m-\ell)} \right)
\\
&\qquad + \hat{\mathbf{R}}^{A}\begin{pmatrix} s,t \\ u, v \end{pmatrix}.
\end{align*}
    
At this stage, we extract the principal contribution by replacing the inner
$\ell$--sum with the corresponding $m$--fold term, and we include the resulting
difference in the remainder:
\begin{align*}
&=
\sum_{i=0}^{\lambda+\kappa-2}
\sum_{m=i}^{\lambda+\kappa-2}
\frac{1}{m!}\,
D^m g\left(Z^{(0,0)}_{t,u}\right)
\circ
\left(
Z^{\otimes(i-1)}_{s,u}\otimes R^1_{s,t}(u)\otimes Z^{\otimes(m-i)}_{t,u}
\right)
\circ
(\Delta^{\boxtimes}_{>0})^{m}
\circ
\left[id \boxtimes L_{\mathbf{Y}^{<\lambda}_{u,v}}\right]
\\
&\qquad + \mathbf{R}^{A}\begin{pmatrix} s,t \\ u, v \end{pmatrix},
\end{align*}
where $\mathbf{R}^{A}\begin{pmatrix} s,t \\ u, v \end{pmatrix}$ absorbs both
$\hat{\mathbf{R}}^{A}\begin{pmatrix} s,t \\ u, v \end{pmatrix}$ and the truncation error arising from the
replacement of the inner $\ell$--sum by
$(\Delta^{\boxtimes}_{>0})^{m}\circ[id \boxtimes L_{\mathbf{Y}^{<\lambda}_{u,v}}]$. The latter term is the second level analogous of the term $C$ in the expansion of the first-level remainder, and can be seen to satisfy inequality \eqref{remainder_remainder_norm}.

We now turn to the terms $B_{s,t}(v)$ and $C_{s,t}(v)$, defined respectively as 
\begin{align*}
B_{s,t}(v)
&:=
\sum_{i = 1}^{\kappa -1} \sum_{m = 0}^{\lambda -1} \frac{1}{(\kappa-1)!} \frac{1}{m!} \binom{\kappa-1}{i} \int_0^1 (1-r)^\kappa D^{\kappa+m}g((1-r)Z^{(0,0)}_{s, v} + rZ^{(0,0)}_{t, v}) dr\\
& \qquad \circ \left((Z_{s,v} \circ [L_{\mathbf{X}^{<\kappa}_{s,t}} \boxtimes id])^{\otimes i + m} \circ \left((\tilde{\Delta} \boxtimes \Delta)^{i} \otimes ( id \boxtimes \tilde{\Delta})^{m}\right) \otimes (Z^{0,0}_{t,v} - Z^{0,0}_{s,v})^{\otimes \kappa-i}\right), \\
C_{s,t}(v) &:= \sum_{i = 1}^{\kappa -1} \sum_{n = i}^{\kappa -1} \sum_{m = 0}^{\lambda -1} \sum_{p=0}^{n-i} \frac{1}{n!} \frac{1}{m!} \binom{n}{i} D^{n+m}g(Z^{(0,0)}_{s,v})\\
& \qquad \circ (Z_{s,v} \circ [L_{\mathbf{X}^{<\kappa}_{s,t}} \otimes id])^{\otimes i + m} \\
&\qquad \circ \left((\tilde{\Delta} \boxtimes \Delta))^{i} \otimes ( id \boxtimes \tilde{\Delta})^{m} \otimes \left(Z_{s,v} \circ \left[\mathbf{X}^{\geq 1 }_{s,t} \boxtimes id \right]\right)^{\otimes p-1} \otimes R^{1, (0,0)}_{s,t}(v) \otimes \left(Z^{(0,0)}_{t,v} - Z^{(0,0)}_{s,v}\right)^{\otimes n-i-p} \right). 
\end{align*}
A procedure analogous to the one used for the term $A$ shows that
\begin{align*}
    &B_{s,t}(v) = B_{s,t}(u) \circ [id \boxtimes L_{\mathbf{X}^{< \kappa}}] + \mathbf{R}^{B}\begin{pmatrix} s,t \\ u, v \end{pmatrix}, \\
    &C_{s,t}(v) = C_{s,t}(u) \circ [id \boxtimes L_{\mathbf{X}^{< \kappa}}] + \mathbf{R}^{C}\begin{pmatrix} s,t \\ u, v \end{pmatrix}. \\
\end{align*}
Where $\mathbf{R}^{B}\begin{pmatrix} s,t \\ u, v \end{pmatrix}$ and $\mathbf{R}^{C}\begin{pmatrix} s,t \\ u, v \end{pmatrix}$
satisfy the bound \eqref{remainder_remainder_norm}.

We finally treat the term $E$. Set
\begin{align*}
E_{s,t}(v)
&:=
\sum_{i = 1}^{\kappa -1} \sum_{n = i}^{\kappa -1} \sum_{m = 0}^{\lambda -1}  \frac{1}{n!} \frac{1}{m!} \binom{n}{i} D^{n+m}g(Z^{(0,0)}_{s,v})\circ Z^{\otimes m + n}_{s,v} \circ \Bigg(  \left[ id \boxtimes \tilde{\Delta}\right]^{\otimes m}   \\
&\qquad \otimes \left(\sum_{i = 1}^{n}\binom{n}{i} [L_{\mathbf{X}^{<\kappa}_{s,t}} \otimes id]^{\otimes i } \circ (\tilde{\Delta} \boxtimes \Delta)^{i}  \otimes \left[\mathbf{X}^{\ge 1}_{s,t} \boxtimes \mathbf{1} \right]^{\otimes n-i} - (\tilde{\Delta} \boxtimes \Delta)^{n}  \otimes [L_{\mathbf{X}^{<\kappa}_{s,t}} \otimes id] \right)\Bigg)\\ 
&= \sum_{i = 1}^{\kappa -1} \sum_{n = i}^{\kappa -1} \sum_{m = 0}^{\lambda -1}  \frac{1}{n!} \frac{1}{m!} \binom{n}{i} D^{n+m}g(Z^{(0,0)}_{s,v})\circ Z^{\otimes m + n}_{s,v} \circ \Bigg(  \left[ id \boxtimes \tilde{\Delta}\right]^{\otimes m} \otimes  G(\mathbf{X}^{\geq 1}_{s,t})\Bigg)
\end{align*}
proceeding in a totally analogous way to what we did above, we obtain expanding $Z$ in the second parameter and grouping the terms that present at least one remainder term $R^2$ in $\hat{\mathbf{R}}^{E}$,
\begin{align*}
&=
\sum_{i=1}^{\kappa-1} \sum_{n=i}^{\kappa-1} \sum_{m=0}^{\lambda-1} 
\frac{1}{n!} \frac{1}{m!} \binom{n}{i}
D^{n+m} g(Z^{(0,0)}_{s,v}) \circ Z^{\otimes m+n}_{s,u} \\
&\quad \circ
\Bigg(
    \left[ id \boxtimes L_{\mathbf{Y}^{<\lambda}_{u,v}} \right]
    \circ \left[ id \boxtimes \tilde{\Delta} \right]^{\otimes m}  \otimes
    \left[ id \boxtimes L_{\mathbf{Y}^{<\lambda}_{u,v}} \right]^{\otimes n}
    \circ G(\mathbf{X}^{\ge 1}_{s,t})
\Bigg) \\
&\quad
+ \hat{\mathbf{R}}^{E} \begin{pmatrix} s,t \\ u,v \end{pmatrix} .
\end{align*}

From definition of $R^2$ it can be seen that the term $\hat{\mathbf{R}}^{E}\begin{pmatrix} s,t \\ u, v \end{pmatrix}$ satisfies \eqref{remainder_remainder_norm}.\newline
Subsequently, by expanding $g$ around $Z^{(0,0)}_{s,u}$ and absorbing the Taylor remainder in $\hat{\mathbf{R}}^{E}\begin{pmatrix} s,t \\ u, v \end{pmatrix}$
\begin{align*}
        = \sum_{i = 1}^{\kappa -1} \sum_{n = i}^{\kappa -1} \sum_{m = 0}^{\lambda -1} \sum_{l=m}^{\lambda -1} &\frac{1}{n!} \frac{1}{l!} \binom{n}{i} \binom{l}{m} D^{n+l}g(Z^{(0,0)}_{s,u})\circ Z^{\otimes l + n}_{s,u}\\ 
        & \circ \Bigg( \left( \left[id \boxtimes L_{\mathbf{Y}^{< \lambda}_{u,v}}\right] \circ \left[ id \boxtimes \tilde{\Delta}\right]\right)^{\otimes m} 
         \otimes  \left(\left[id \boxtimes L_{\mathbf{Y}^{< \lambda }_{u,v}}\right]\right)^{\otimes n} \circ G(\mathbf{X}^{\geq 1}_{s,t})\otimes \left[ \mathbf{1} \boxtimes \mathbf{Y}^{\geq 1}_{u,v} \right]^{\otimes l-m}\Bigg) \\
         \qquad + \hat{\mathbf{R}}^{E}\begin{pmatrix} s,t \\ u, v \end{pmatrix}.
\end{align*}
Again, the newly absorbed term satisfies \eqref{remainder_remainder_norm}.\newline
Finally proceeding as we did in the previous proof in order to recover $C$ we obtain the desired identity i.e.
\begin{align*}
E_{s,t}(v) = E_{s,t}(u)\circ [id \boxtimes L_{\mathbf{Y}^{< \lambda}_{u,v}}] + \mathbf{R}^{E}\begin{pmatrix} s,t \\ u, v \end{pmatrix}, 
\end{align*}
where $\mathbf{R}^{E}\begin{pmatrix} s,t \\ u, v \end{pmatrix}$  is the sum of this newly found second level remainder term and $\hat{\mathbf{R}}^{E}\begin{pmatrix} s,t \\ u, v \end{pmatrix}$. This allows to conclude by noting that $\mathbf{R}^{E}\begin{pmatrix} s,t \\ u, v \end{pmatrix}$ respects the bound \eqref{remainder_remainder_norm} and that, from construction, there exists a $K > 0$ such that 
\[
\big\| \mathbf{R}^{F, (j,k)} \big\|_{(\frac{p}{\kappa - j},\frac{q}{\lambda - k}), (\omega_\mathbf{X}, \omega_\mathbf{Y})} \leq K \| Z \|_{\mathcal{D}_{\mathbf{X}, \mathbf{Y}}},
\] 
for all $j \in [0:\kappa)$, $k \in [0:\lambda)$, $F \in \{A, B, C, E\}$.
\end{proof}

\end{document}
